\chapter{Risque en temps réel}\label{ch:pl:real-time-risk}

Chacun des deux desks est resté dans sa limite tout l'après-midi. Ensemble, ils se sont retrouvés vendeurs à près du double de l'appétit de risque de la société sur le même indice, et
personne ne le voit avant le \omterm{def:pl:the-platform-map:batch}{batch de fin de journée}, parce que rien n'additionne les desks pendant qu'ils négocient. Le système de chaque desk répond à la question
qui lui est posée --- suis-je dans ma limite~? --- et la question que personne n'a posée à aucun système est celle qui compte. Ce chapitre construit
l'agrégateur qui la pose à chaque exécution et à chaque mouvement de marché~: les expositions sont additionnées le long de la hiérarchie de risque de manière incrémentale, les sensibilités sont mises en cache et
rafraîchies lorsque le marché bouge, les limites contrôlées dans la boucle, et un \omterm{def:pl:real-time-risk:whatif}{contrôle what-if} avant l'envoi d'un ordre.

\section{Ce que signifie le temps réel pour le risque}

\begin{definition}[Système de risque en temps réel]\label{def:pl:real-time-risk:system}
Un \emph{système de risque en temps réel}\index{système de risque en temps réel} calcule en continu les expositions de la société et l'utilisation de ses limites à partir des
exécutions du \omterm{def:pl:position-and-pnl-services:service}{service de positions} et des mouvements du marché, assez vite pour qu'un dépassement soit repéré pendant que les opérations de trading qui le causent peuvent encore être arrêtées, et
les restitue avec leurs horodatages.
\end{definition}

Le moteur de risque de fin de journée du livre~6 (chapitre~29) réévalue chaque position selon des centaines de scénarios une fois par nuit~; la passerelle pré-trade
du livre~13 (chapitre~22) contrôle un ordre en microsecondes par rapport aux limites qu'elle tient en mémoire. Le risque en temps réel se situe entre les deux~: la durée se mesure en secondes, pas en
nanosecondes ni en heures, et englobe toute la société, et non un seul ordre. Ses chiffres sont moins fins que ceux du moteur nocturne (sensibilités, et non
réévaluation complète) et son périmètre plus large que celui de la passerelle (chaque desk à la fois). La rapidité de mise à disposition des données de risque est l'un des principes du BCBS~239
(livre~6, chapitre~29)~: un chiffre disponible le lendemain matin décrit un risque déjà pris.

\section{Agrégation le long de la hiérarchie}

\begin{definition}[Agrégation des expositions]\label{def:pl:real-time-risk:aggregation}
L'\emph{agrégation des expositions}\index{agrégation des expositions} additionne les expositions de chaque transaction --- ici des dollars de delta sur l'indice --- le long de la
hiérarchie de risque (transaction, book, desk, société) de sorte que chaque nœud porte la somme de ses descendants et puisse être comparé à sa propre limite.
\end{definition}

La limite d'un desk et celle de la société ne forment pas une partition~: l'appétit de la société est en général inférieur à la somme des limites de ses desks, parce que ceux-ci
ne sont pas censés utiliser leurs limites au même moment et dans le même sens. C'est exactement ce qu'ils font parfois. La hiérarchie du chapitre
compte une société, deux desks et un book chacun~; l'appétit de la société et la limite de chaque desk valent 50~millions de dollars de delta sur l'indice, ce que
la règle de \texttt{firm.riskctl} autorise (la limite d'un desk ne dépasse jamais celle de la société).

\begin{omfigure}
\begin{tikzpicture}[font=\scriptsize,box/.style={draw,rounded corners=2pt,align=center,minimum width=3.4cm,minimum height=0.9cm},>=Stealth]
\node[box,fill=omThm!15] (f) at (0,0) {société~: limite 50~M\$\\ clôture $-\$90.8$~M\$ (limite dépassée)};
\node[box,fill=omDef!12] (a) at (-3,-1.5) {desk A~: limite 50~M\$\\ clôture $-\$47.6$~M\$};
\node[box,fill=omDef!12] (b) at (3,-1.5) {desk B~: limite 50~M\$\\ clôture $-\$43.2$~M\$};
\node[box,fill=gray!10] (a1) at (-3,-3) {book~: puts sur l'indice\\ 21\,000 unités, six lignes};
\node[box,fill=gray!10] (b1) at (3,-3) {book~: future sur l'indice\\ 8\,800 à découvert};
\draw[->] (a1) -- (a); \draw[->] (b1) -- (b); \draw[->] (a) -- (f); \draw[->] (b) -- (f);
\node[font=\scriptsize\itshape,align=left,anchor=west] at (5.0,-2.2) {chaque exécution et chaque tick remonte\\ sa variation vers le haut};
\end{tikzpicture}

\omcaption{La hiérarchie de risque du chapitre à la clôture, avec la limite et le delta de chaque nœud. Chaque nœud est dans sa limite, sauf celui
qu'aucun système de desk ne calcule.}\label{fig:pl:real-time-risk:hierarchy}
\end{omfigure}

\section{Mises à jour incrémentales}

\begin{definition}[Agrégation incrémentale]\label{def:pl:real-time-risk:incremental}
L'\emph{agrégation incrémentale}\index{agrégation incrémentale} met à jour les agrégats en n'ajoutant que la variation provoquée par chaque événement --- l'exposition d'une exécution,
la variation d'une transaction après un mouvement de marché --- au nœud et à ses ascendants, au lieu de recalculer l'arbre à partir de toutes les transactions.
\end{definition}

\begin{definition}[Cache de sensibilités]\label{def:pl:real-time-risk:cache}
Un \emph{cache de sensibilités}\index{cache de sensibilités} conserve les sensibilités de chaque transaction (delta, gamma, vega) calculées par la bibliothèque de pricing à
un même état du marché, s'en sert pour reporter les expositions au fil des mouvements du marché, et les recalcule lorsque le marché a bougé de plus d'un
seuil depuis ce calcul.
\end{definition}

L'agrégateur du chapitre (\texttt{firm.rtrisk}) conserve, pour chaque transaction, les dollars de delta qu'elle apporte actuellement. Une exécution ajoute la nouvelle
contribution moins l'ancienne au book de la transaction et à chaque ascendant. Un tick de marché fait avancer chaque delta en cache avec son gamma, $(\Delta +
\Gamma\,\mathrm{d}S)\,q\,S$, et applique les écarts~; lorsque le spot a bougé de plus de 0.5~\% depuis le remplissage du cache, chaque
instrument est réévalué par la bibliothèque du livre~5 (\texttt{firm.pricing}, bump-and-reprice), puis les écarts sont de nouveau appliqués. Les agrégats ne
sont jamais reconstruits~; les tests vérifient qu'ils égalent toujours une reconstruction complète.

\omcode{code/firm/rtrisk/firm_rtrisk.py}{50}{76}{Le cache de sensibilités~: grecques issues de la bibliothèque de pricing à un spot, périmées au-delà
d'un seuil.}\label{lst:pl:real-time-risk:cache}

\omcode{code/firm/rtrisk/firm_rtrisk.py}{99}{134}{Agrégation incrémentale~: la contribution d'une transaction, sa variation ajoutée à chaque ascendant, et
limites contrôlées à chaque exécution et à chaque tick.}\label{lst:pl:real-time-risk:aggregator}

\begin{definition}[Instantané de risque]\label{def:pl:real-time-risk:snapshot}
Un \emph{instantané de risque}\index{instantané de risque} est l'état de l'agrégateur à un instant donné~: l'exposition de chaque nœud, l'heure de marché qu'il reflète, le
moment et le spot des sensibilités utilisées, leur ancienneté et les limites dépassées.
\end{definition}

Un \omterm{def:pl:real-time-risk:snapshot}{instantané de risque} qui ne dit pas depuis quand ses sensibilités ont été calculées ne peut pas être tenu pour fiable~; un instantané qui le dit peut servir avec discernement. À la clôture de
l'après-midi du chapitre, le cache avait été rafraîchi pour la dernière fois 65 minutes plus tôt (3\,890~secondes)~: le spot n'avait pas bougé de 0.5~\% depuis, de sorte que les
deltas reportés restaient dans la précision du seuil.

\section{Sensibilités, réévaluation complète et écart}

Reporter les expositions à partir de sensibilités en cache est rapide et approximatif~; la réévaluation complète est exacte et lente. L'écart se mesure en valorisant
les options du desk A à la clôture selon les deux méthodes, pour des mouvements allant de $-10\%$ à $+10\%$ (la \cref{fig:pl:real-time-risk:gap})~: le P\&L delta-gamma issu de
sensibilités calculées à la clôture face à la variation de la valeur complète des options.

\begin{omfigure}
\begin{tikzpicture}
\begin{axis}[width=11.5cm,height=5.6cm,xlabel={mouvement de l'indice (\%)},ylabel={erreur du P\&L delta-gamma (\% de la valeur complète)},xmin=-10,xmax=10,
  ymin=0,ymax=25,grid=major,grid style={gray!20},tick label style={font=\small},label style={font=\small},table/col sep=comma]
\addplot[color=omDef,very thick] table[x=move_pct,y=error_pct]{figdata/platforms/18-real-time-risk/approximation.csv};
\draw[omThm,dashed] (axis cs:-10,1) -- (axis cs:10,1) node[pos=0.5,above,font=\scriptsize]{1~\%};
\draw[black,dotted] (axis cs:3.5,0) -- (axis cs:3.5,25) node[pos=0.8,right,font=\scriptsize]{$+3.5~\%$};
\draw[black,dotted] (axis cs:-6.75,0) -- (axis cs:-6.75,25) node[pos=0.8,left,font=\scriptsize]{$-6.75~\%$};
\end{axis}
\end{tikzpicture}

\omcaption{L'erreur du P\&L calculé à partir du delta et du gamma en cache, en proportion de celui de la réévaluation complète, pour les 21\,000 puts sur l'indice
du desk A à la clôture, selon l'ampleur du mouvement de l'indice. Elle est nulle en $\pm1~\%$ (là où portaient les bumps des grecques), inférieure à 1~\% entre $-6.75~\%$ et $+3.5~\%$,
et croît vite à la hausse, là où la convexité des puts change le plus. Données~: \texttt{fig\_rtrisk.py}.}\label{fig:pl:real-time-risk:gap}
\end{omfigure}

Dans la plage de mouvements qu'un système en temps réel rencontre en une journée, l'approximation est excellente~: l'erreur est exactement nulle en $\pm 1~\%$, parce que la bibliothèque
a calculé le delta et le gamma à partir de bumps de $\pm 1~\%$ et la formule quadratique passe par eux, et reste sous 1~\% de la réévaluation complète entre
$-6.75~\%$ et $+3.5~\%$. Au-delà, elle se dégrade vite --- 3.1~\% à $+5~\%$, 23.6~\% à $+10~\%$ --- ce qui explique que le cache soit rafraîchi pour des mouvements bien
plus petits que ceux-là, et pourquoi les chiffres de stress viennent de la réévaluation complète (livre~6, chapitre~29), et non des sensibilités.

\section{Limites dans la boucle}

\begin{definition}[Contrôle what-if]\label{def:pl:real-time-risk:whatif}
Un \emph{contrôle what-if}\index{contrôle what-if} calcule ce qu'une transaction proposée ferait à l'exposition de chaque nœud au-dessus de son livre,
ainsi que les limites qu'elle franchirait, avant l'envoi de la transaction, sans modifier l'état de l'agrégateur.
\end{definition}

\omcode{code/firm/rtrisk/firm_rtrisk.py}{144}{148}{Le contrôle what-if~: la contribution de la transaction, ajoutée à ses ascendants sur une copie, et les
limites qu'elle franchirait.}\label{lst:pl:real-time-risk:whatif}

Un \omterm{def:pl:real-time-risk:whatif}{contrôle what-if} est la contribution du système en temps réel au contrôle pré-trade~: le filtre du livre~13 vérifie un ordre face aux limites en
microsecondes, le what-if lui indique si la limite au niveau de la société a encore de la marge. À la clôture de l'après-midi du chapitre, le desk~B qui demande à vendre 400
contrats à terme supplémentaires serait informé que son desk reste à $-45.2$~millions de dollars, dans sa limite, et que la société passe à $-92.8$~millions de dollars, très au-delà.

\section{Deux desks, un seul indice}

L'après-midi du chapitre court de 13~h~00 à 17~h~30 sur un indice qui part de 5\,000, avec des mouvements simulés à graine fixée (volatilité annuelle de 20~\% et dérive de
$-1.5~\%$ sur l'après-midi~; il clôt à 4\,908.5). Le desk~A achète des puts sur indice --- six lignes cotées, trois strikes et deux échéances, 1\,000 unités
par lot --- et le desk~B vend des contrats à terme sur indice, 400 par lot, chacun à des moments aléatoires, environ toutes les cinq minutes. Avant chaque transaction, le desk
vérifie sa propre exposition et ne négocie que s'il reste sous 90~\% de sa limite. Les exécutions passent par le \omterm{def:pl:position-and-pnl-services:service}{service de positions} du chapitre~17~;
l'agrégateur de la société additionne tout.

\begin{omfigure}
\begin{tikzpicture}
\begin{axis}[width=12cm,height=6cm,xlabel={heure de la journée (heures)},ylabel={delta (millions de \$)},xmin=13,xmax=17.5,ymin=-100,ymax=5,
  grid=major,grid style={gray!20},tick label style={font=\small},label style={font=\small},table/col sep=comma,
  legend style={font=\footnotesize,at={(0.5,1.03)},anchor=south,legend columns=3}]
\addplot[color=omThm,very thick,const plot] table[x=hour,y=firm_m]{figdata/platforms/18-real-time-risk/exposure.csv};
\addlegendentry{société}
\addplot[color=omDef,thick,const plot] table[x=hour,y=desk_a_m]{figdata/platforms/18-real-time-risk/exposure.csv};
\addlegendentry{desk A (puts)}
\addplot[color=omMeth,thick,dashed,const plot] table[x=hour,y=desk_b_m]{figdata/platforms/18-real-time-risk/exposure.csv};
\addlegendentry{desk B (contrats à terme)}
\draw[black,thick] (axis cs:13,-50) -- (axis cs:17.5,-50) node[pos=0.72,below,font=\scriptsize]{chaque limite~: $-50$~millions de dollars};
\draw[omThm,dotted,thick] (axis cs:14.397,-100) -- (axis cs:14.397,5) node[pos=0.93,right,font=\scriptsize]{dépassement société 14:23:50};
\end{axis}
\end{tikzpicture}

\omcaption{Delta de l'indice pour la société et pour chaque desk tout au long de l'après-midi, dans l'agrégateur. Aucun des deux desks ne dépasse sa limite --- la plus forte exposition du desk~A
est de $-49.3$~millions de dollars, celle du desk~B de $-43.6$~millions de dollars --- alors que la société franchit la sienne à 14:23:50 et termine à $-90.8$~millions de dollars.
Données~: \texttt{fig\_rtrisk.py}.}\label{fig:pl:real-time-risk:exposure}
\end{omfigure}

\begin{example}[Deux desks, un seul indice]\label{ex:pl:real-time-risk:desks}
Les deux desks ont fait ce que leurs systèmes permettaient~: le desk~A a réalisé 21 transactions et détient 21\,000 puts, le desk~B en a réalisé 22 et est vendeur de 8\,800 contrats à terme~; aucun des deux
n'a dépassé sa limite à aucun moment. La limite de la société a été franchie à 14:23:50, seule alerte de l'agrégateur de l'après-midi~; à la clôture, la
société affiche $-90.8$~millions de dollars de delta, soit 1.82 fois son appétit pour le risque (desk~A $-47.6$~millions de dollars, desk~B $-43.2$~millions de dollars). Sans l'agrégateur, le
\omterm{def:pl:the-platform-map:batch}{batch de fin de journée} l'aurait détecté trois heures après le franchissement. Le cache a été rafraîchi six fois~; l'agrégateur a effectué 10\,570 mises à jour.
\end{example}

\begin{remark}[Qui agit sur l'alerte]\label{rem:pl:real-time-risk:who}
Un franchissement au niveau de la société n'appartient à aucun desk. L'alerte de l'agrégateur doit atteindre quelqu'un qui a autorité sur les deux --- le responsable des risques qui peut
bloquer les nouveaux ordres qui augmentent le risque via le filtre, ou réduire sur-le-champ la limite d'un desk (le \omterm{def:pl:signal-serving-and-parameter-management:store}{magasin de paramètres} du chapitre~16, avec ses approbations) ---
et le cadre des limites doit dire à l'avance dont le risque sera réduit en premier. Un chiffre en temps réel que personne ne peut exploiter n'est qu'un rapport plus rapide.
\end{remark}

\section{Tutoriel~: deux desks et un agrégateur}\label{tut:pl:real-time-risk:tut}

\begin{tutorial}
\textbf{Objectif.} Agréger incrémentalement les expositions des deux desks, repérer le franchissement au niveau de la société que les desks ne pouvaient pas voir, et mesurer l'erreur d'approximation du
\omterm{def:pl:real-time-risk:cache}{cache de sensibilités}. \textbf{État final~:} \cref{fig:pl:real-time-risk:exposure,fig:pl:real-time-risk:gap}.
\begin{tutsteps}
\item \textbf{L'après-midi}~: \texttt{pl\_rtrisk.afternoon()}~: la trajectoire de l'indice, les transactions des desks, les alertes de l'agrégateur.
\item \textbf{Instantanés}~: \texttt{agg.snapshot(t)} à tout moment~: expositions, âge du cache, franchissements.
\item \textbf{\omterm{def:pl:real-time-risk:whatif}{Contrôle what-if}}~: \texttt{what\_if\_at\_close(r)}.
\item \textbf{L'écart}~: \texttt{approximation\_error(r)} et \texttt{one\_percent\_move}.
\end{tutsteps}
\textbf{Pour aller plus loin.} Abaisser le seuil de rafraîchissement à 0.1~\% et compter les réévaluations~; donner au desk~A des calls au lieu de puts et voir si
la société franchit encore sa limite.
\end{tutorial}

\section{Construction~: l'agrégateur en temps réel}\label{bld:pl:real-time-risk:rtrisk}

\begin{build}
\bfield{Rôle} Les expositions de la société et la consommation des limites, chaque nœud de la hiérarchie, mises à jour à chaque exécution et à chaque tick, avec les réponses what-if pour
les transactions proposées.
\bfield{Interface} \texttt{Hierarchy}~; \texttt{SensitivityCache(md, underlying, threshold)}~;
\texttt{Aggregator(hierarchy, cache, limits)} avec \texttt{on\_fill}, \texttt{on\_tick}, \texttt{exposure}, \texttt{snapshot}, \texttt{what\_if}, \texttt{alerts}~;
\texttt{Future}~; \texttt{delta\_gamma\_pnl}, \texttt{full\_pnl}.
\bfield{Règles} Les agrégats sont mis à jour par différences, jamais reconstruits~; les sensibilités viennent de la bibliothèque de valorisation et sont rafraîchies au-delà
d'un seuil~; chaque instantané porte l'âge de ses sensibilités~; les limites sont vérifiées à chaque mise à jour~; un what-if ne modifie jamais l'état.
\bfield{Tests d'acceptation} \texttt{code/firm/rtrisk/tests/}~: la hiérarchie et les sommes incrémentales face à une reconstruction complète~; les ticks dans le
seuil (porté avec le gamma) et au-delà (réévalué)~; limites, alertes et what-if~; le P\&L delta-gamma face à la réévaluation complète pour un petit
mouvement.
\bfield{Pour aller plus loin} Expositions au vega et aux taux sur la même hiérarchie~; de nombreux sous-jacents avec un chiffre consolidé de la société tenant compte des corrélations~; l'agrégateur alimenté
à partir du journal d'événements du chapitre~24.
\end{build}

\begin{omsources}
\item One Quant Book~5, chapitres 4 et 28 (grecques, bibliothèque de valorisation)~; One Quant Book~6, chapitres 21 et 29 (sensibilités, le moteur de
risque et sa hiérarchie, BCBS~239)~; One Quant Book~11, chapitre~27, et livre~13, chapitre~22 (limites et filtre pré-trade).
\end{omsources}

\section{Exercices}

\begin{exercise}[$\star$]\label{exo:pl:real-time-risk:1}
Pour un put de delta $-0.4$ sur 1\,000 unités d'un indice à 5\,000, calculer son delta en dollars et déterminer combien de transactions de ce type tiennent sous 90~\% de la
limite de 50~millions de dollars~?
\end{exercise}

\begin{exercise}[$\star$]\label{exo:pl:real-time-risk:2}
Expliquer pourquoi l'erreur d'approximation est exactement nulle en $\pm 1\%$ dans la \cref{fig:pl:real-time-risk:gap}~?
\end{exercise}

\begin{exercise}[$\star$]\label{exo:pl:real-time-risk:3}
Dire ce que le \omterm{def:pl:real-time-risk:whatif}{contrôle what-if} indique au desk~B à la clôture, et expliquer pourquoi le système propre au desk~B ne peut pas le lui dire~?
\end{exercise}

\begin{exercise}[$\star\star$]\label{exo:pl:real-time-risk:4}
Expliquer pourquoi l'erreur croît plus vite à la hausse qu'à la baisse pour un livre de puts longs~?
\end{exercise}

\begin{exercise}[$\star\star$]\label{exo:pl:real-time-risk:5}
L'instantané de la clôture utilisait des sensibilités vieilles de 65 minutes. Quand cela est-il acceptable, et qu'est-ce qui le rendrait inacceptable~?
\end{exercise}

\begin{exercise}[$\star\star$]\label{exo:pl:real-time-risk:6}
Pourquoi l'\omterm{def:pl:real-time-risk:incremental}{agrégation incrémentale} est-elle vérifiée contre une reconstruction dans les tests, et quel défaut le contrôle permettrait-il de détecter~?
\end{exercise}

\begin{exercise}[$\star\star\star$]\label{exo:pl:real-time-risk:7}
\emph{Programmation.} Exécuter \texttt{afternoon} avec un seuil de rafraîchissement de 0.1~\% puis de 2~\%. Combien de
réévaluations chaque seuil entraîne-t-il, et l'exposition de la société à la clôture change-t-elle~?
\end{exercise}

\begin{exercise}[$\star\star\star$]\label{exo:pl:real-time-risk:8}
\emph{Trouver la faille.} «~Chaque desk dispose d'une limite et d'un contrôle pré-trade, donc la société ne peut pas dépasser la somme des limites des desks, et nous avons fixé
cette somme en dessous de l'appétit pour le risque de la société.~»
\end{exercise}

\section{Problème~: Deux desks, un même indice}

\begin{problem}[{Problème du week-end --- une violation qu'aucun système ne pouvait voir}]\label{pb:pl:real-time-risk:1}
L'après-midi du chapitre, les desks et l'agrégateur.

\medskip\noindent\textbf{Partie I --- Le système.}
\begin{enumerate}
\item Que signifie le temps réel ici, par rapport au moteur de risque nocturne et à la barrière pré-trade~?
\item Comment les expositions sont-elles agrégées, et pourquoi les limites des desks ne forment-elles pas une partition de celle de la société~?
\item Qu'apporte une mise à jour incrémentale, sur une exécution et sur un tick~?
\item Quand le cache est-il rafraîchi, et par quoi~?
\item Que contient un instantané~?
\end{enumerate}

\medskip\noindent\textbf{Partie II --- L'après-midi.}
\begin{enumerate}[resume]
\item Qu'a négocié chaque desk, et qu'a-t-il contrôlé~?
\item Quelles étaient les plus fortes expositions des desks~?
\item Quand la limite de la société a-t-elle été dépassée~?
\item Quelle était l'exposition de la société à la clôture, en multiple de son appétit~?
\item Que dit le \omterm{def:pl:real-time-risk:whatif}{contrôle what-if} à la clôture~?
\end{enumerate}

\medskip\noindent\textbf{Partie III --- L'approximation.}
\begin{enumerate}[resume]
\item Comment le P\&L du cache se compare-t-il à la réévaluation complète~?
\item Entre quels mouvements l'erreur est-elle inférieure à 1~\%~?
\item Quelle est-elle à $+5~\%$ et à $+10~\%$~?
\item Pourquoi est-elle nulle à $\pm1~\%$~?
\item Quelles conséquences pour les chiffres de stress~?
\end{enumerate}

\medskip\noindent\textbf{Partie IV --- Le verdict.}
\begin{enumerate}[resume]
\item Énoncer le \emph{résultat nommé}~: l'heure à laquelle la limite au niveau de la société a été dépassée alors que chaque desk restait dans la sienne, l'exposition à
la clôture, et la taille du mouvement de marché au-delà duquel l'erreur de P\&L du cache dépasse un pour cent de la réévaluation complète.
\item Qui doit recevoir l'alerte, et que doit-on pouvoir faire~?
\item Comment fixer le seuil de rafraîchissement~?
\item Que faudrait-il agréger d'autre en temps réel~?
\item En une phrase~: pourquoi faut-il additionner le risque au fur et à mesure qu'il est pris~?
\end{enumerate}
\end{problem}

\section{Questions d'entretien}

\begin{interviewq}[$\star$ \iqroles{développeur, chercheur}]\label{iq:pl:real-time-risk:1}
Deux desks restent chacun dans leurs limites. Comment la société peut-elle sortir des siennes~?
\end{interviewq}

\begin{interviewq}[$\star\star$ \iqroles{développeur}]\label{iq:pl:real-time-risk:2}
Comment maintiendriez-vous à jour les expositions de l'ensemble de la société alors que des milliers d'exécutions et de ticks arrivent chaque seconde~?
\end{interviewq}

\begin{interviewq}[$\star\star$ \iqroles{chercheur}]\label{iq:pl:real-time-risk:3}
Quand une approximation delta--gamma du P\&L est-elle suffisante, et quand devez-vous réévaluer~?
\end{interviewq}

\begin{interviewq}[$\star\star$ \iqroles{développeur}]\label{iq:pl:real-time-risk:4}
Qu'est-ce qu'un \omterm{def:pl:real-time-risk:snapshot}{instantané de risque} doit dire à son lecteur, au-delà des chiffres~?
\end{interviewq}

\begin{interviewq}[$\star\star\star$ \iqroles{développeur}]\label{iq:pl:real-time-risk:5}
Concevez un \omterm{def:pl:real-time-risk:system}{système de risque en temps réel} pour une société de trading à plusieurs desks.
\end{interviewq}
